\documentclass[12pt,a4paper]{article}

\usepackage[spanish]{babel}
\usepackage[utf8]{inputenc}
\usepackage[T1]{fontenc}
\usepackage{amsmath}
\usepackage{amsfonts}
\usepackage{geometry}
\usepackage{xcolor}
\usepackage{hyperref}
\usepackage{booktabs}

\geometry{margin=2.5cm}
\hypersetup{colorlinks=true, urlcolor=blue, linkcolor=black, citecolor=blue}

\title{\textbf{Informe de Investigación: Plataforma LogiLearn de Lógica Simbólica}}
\author{
  Equipos Sinergia, Los Hijos de Linus, Modus Innova y Linus \\
  Curso Lógica Simbólica (MATG1001) \\
  Docente: Dr. Mardo Victor Gonzales Herrera \\
  Universidad Nacional Pedro Ruiz Gallo \\
  Escuela Profesional de Ingeniería de Sistemas \\
  Semestre 2026 I $\cdot$ II Ciclo
}
\date{Agosto de 2026}

\begin{document}

\maketitle

\section*{Resumen}
El presente informe documenta el diseño, desarrollo, validación y resultados de \textbf{LogiLearn}, una
plataforma web colaborativa construida como producto del curso de Lógica Simbólica (MATG1001). La
plataforma implementa funcionalmente el 100\% de los temas del sílabo mediante \textit{motores reales
de evaluación lógica}: tablas de verdad, inferencias, cuantificadores y teoría de conjuntos. Se
describe el marco tecnológico (Vue 3, Vite, Tailwind CSS v4, TypeScript, Vitest), la metodología de
trabajo en cuatro equipos de dominio con ramificaciones Git y Pull Requests, la arquitectura por
módulo, la estrategia de pruebas (175 pruebas funcionales), un manual de uso y la reflexión sobre los
aprendizajes alcanzados respecto a los resultados de aprendizaje D1, D2 y D3 del sílabo. Los
resultados evidencian que la construcción de software verificable refuerza la comprensión teórica de
la lógica simbólica.

\textbf{Palabras clave:} lógica simbólica, lógica proposicional, cuantificadores, teoría de conjuntos,
Vue, educación interactiva, ingeniería de software educativa.

\newpage

\section{Introducción}
La lógica simbólica constituye el fundamento del razonamiento formal en ingeniería de sistemas. El
curso MATG1001 plantea como resultado de aprendizaje el desarrollo de la capacidad de ``plantear
estrategias de solución a problemas de su entorno, usando el razonamiento lógico y analítico en
diversos contextos'' \cite{silabo}. Sin embargo, la apropiación de conceptos como la validez de
inferencias o el alcance de los cuantificadores suele ser abstracta para el estudiante cuando se
aborda únicamente desde la teoría.

LogiLearn nace para cerrar esa brecha: una plataforma donde cada tema del sílabo es un motor
computacional verificable. El objetivo de este informe es (a) documentar qué se construyó, (b) cómo
se construyó y (c) qué aprendizajes se derivaron, sirviendo a la vez de memoria académica y de manual
de uso.

\subsection{Resultados de aprendizaje del sílabo}
\begin{itemize}
  \item \textbf{D1:} Identifica y aplica las diversas definiciones, teorías y conceptos de la lógica proposicional.
  \item \textbf{D2:} Interpreta y aplica las definiciones, teorías y conceptos de la inferencia lógica.
  \item \textbf{D3:} Discute y analiza los conceptos de la teoría de conjuntos y lo aplica a la lógica simbólica.
\end{itemize}

\subsection{Formulación del problema}
Los estudiantes enfrentan tres barreras: (i) dificultad para verificar manualmente tablas de verdad
de más de tres variables; (ii) ausencia de retroalimentación inmediata al aplicar reglas de
inferencia; y (iii) falta de herramientas para practicar cuantificadores y conjuntos en español.
LogiLearn aborda las tres con motores automatizados.

\section{Marco Tecnológico}
La plataforma se construyó sobre un stack moderno, tipado y orientado a componentes:

\begin{center}
\begin{tabular}{lll}
\toprule
Tecnología & Rol & Versión \\
\midrule
Vue 3 & Framework de componentes (SFC \texttt{.vue}) & 3.5+ \\
Vite & Bundler y servidor de desarrollo & 6.x \\
TypeScript & Tipado estático & 5.x \\
Tailwind CSS v4 & Estilos utilitarios & 4.x \\
Vue Router & Enrutamiento & 4.x \\
Vitest & Pruebas unitarias & 4.x \\
\bottomrule
\end{tabular}
\end{center}

El proyecto se despliega mediante \textbf{GitHub Pages} (producción) y \textbf{Netlify} (vistas
previas de PR), lo que garantiza integración continua y retroalimentación temprana. El tipado
estático (TypeScript) fue decisivo para reducir errores en los motores lógicos.

\section{Marco Teórico}
\subsection{Lógica proposicional}
Una proposición es una afirmación con valor de verdad V o F. Los conectivos son negación ($\neg$),
conjunción ($\wedge$), disyunción ($\vee$), implicación ($\to$) y bicondicional ($\leftrightarrow$).
Una fórmula es \textbf{tautología} si es V para toda asignación, \textbf{contradicción} si es F para
toda asignación, y \textbf{contingencia} en otro caso.

\subsection{Inferencia lógica}
Una inferencia es válida si, siendo verdaderas las premisas, la conclusión es necesariamente verdadera
(regla de separación o \textit{Modus Ponens}: $p,\; p\to q \vdash q$). Otras reglas: \textit{Modus
Tollens} ($\neg q,\; p\to q \vdash \neg p$) y \textit{Silogismo Hipotético} ($p\to q,\; q\to r \vdash
p\to r$).

\subsection{Cuantificadores}
El cuantificador universal $\forall x\, P(x)$ es verdadero si $P(x)$ lo es para todo $x$ en el
dominio $D$; el existencial $\exists x\, P(x)$ lo es si hay al menos un testigo. Leyes de De Morgan
para cuantificadores:
\[
  \neg(\forall x\, P(x)) \equiv \exists x\, \neg P(x), \qquad
  \neg(\exists x\, P(x)) \equiv \forall x\, \neg P(x).
\]

\subsection{Teoría de conjuntos}
Operaciones básicas: unión ($A\cup B$), intersección ($A\cap B$), diferencia ($A-B$), complemento
($A^c$) y diferencia simétrica ($A\triangle B$). La teoría de conjuntos se vincula con la lógica
mediante las leyes de De Morgan conjuntistas: $(A\cup B)^c = A^c\cap B^c$ y $(A\cap B)^c = A^c\cup B^c$.

\section{Metodología}
El trabajo se organizó en \textbf{cuatro equipos de dominio}, cada uno responsable de un módulo del
sílabo:

\begin{center}
\begin{tabular}{llll}
\toprule
Equipo & Sublíder & Tema & Carpeta \\
\midrule
Sinergia & Alexa & Tablas de Verdad & \texttt{src/pages/tablas} \\
Los Hijos de Linus & Arom & Inferencias & \texttt{src/pages/inferencias} \\
Modus Innova & Cristian & Cuantificadores & \texttt{src/pages/cuantificadores} \\
Linus & Jordy & Conjuntos & \texttt{src/pages/conjuntos} \\
\bottomrule
\end{tabular}
\end{center}

\subsection{Ramificaciones y flujo Git}
Se adoptó un modelo de ramas por funcionalidad. Cada equipo desarrolló en su rama y consolidó
mediante \textbf{Pull Requests (PR)} revisados por el líder de proyecto. La integración final se
realizó en la rama \texttt{feat/integrate-sinergia-modus} (PR \#16), donde se unificaron motores,
branding y pruebas. Las ramas de trabajo previas (PR \#8, \#14, \#15) se cerraron al integrarse
manualmente.

\subsection{Calidad y pruebas}
Cada motor cuenta con pruebas unitarias (Vitest). La suite final suma \textbf{175 pruebas
funcionales} que validan parsers, evaluadores, leyes lógicas y componentes. El build de producción se
verifica con \texttt{vue-tsc} + \texttt{vite build}. Distribución de pruebas por módulo:

\begin{center}
\begin{tabular}{lr}
\toprule
Módulo & Pruebas \\
\midrule
Tablas de verdad & 39 \\
Inferencias & 24 \\
Cuantificadores & 30 \\
Conjuntos & 14 \\
Solver/parser & 22 \\
UI / integración & 46 \\
\bottomrule
\end{tabular}
\end{center}

\subsection{Gestión de riesgos}
\begin{itemize}
  \item \textit{Riesgo:} discrepancias de formato entre equipos $\to$ mitigado con componentes UI compartidos.
  \item \textit{Riesgo:} errores en motores lógicos $\to$ mitigado con pruebas unitarias por cada ley.
  \item \textit{Riesgo:} pérdida de avances $\to$ mitigado con PRs y revisión centralizada.
\end{itemize}

\section{Arquitectura y Desarrollo}
La lógica se aisló en \texttt{src/lib/} (motores compartidos) y la interfaz en \texttt{src/pages/} y
\texttt{src/components/ui/} (Button, Card, Badge, ToggleSwitch, OptionPill). A continuación se
describen los módulos.

\subsection{Tablas de Verdad (Equipo Sinergia)}
Motor que parsea proposiciones, recolecta variables, genera combinaciones y evalúa subexpresiones.
\textbf{D1.} Clasifica cada fórmula como tautología, contradicción o contingencia. Incluye 30+
ejercicios (identificar, clasificar, leyes, quiz) y verificación en vivo.

\subsection{Inferencias (Equipo Los Hijos de Linus)}
Demostrador con trazabilidad y diagnóstico. Valida reglas como Modus Ponens, Modus Tollens y
Silogismo Hipotético, mostrando cada paso. \textbf{D2.} El estudiante interpreta y aplica la
inferencia lógica con retroalimentación inmediata.

\subsection{Cuantificadores (Equipo Modus Innova)}
Evalúa $\forall$ y $\exists$ sobre dominios finitos, con rangos encadenados (\texttt{0 < x < 90}),
predicados libres (\texttt{x \% 2 === 0}), negación De Morgan profunda y un \textbf{resolutor paso a
paso} (Bicondicional $\to$ Implicación $\to$ De Morgan $\to$ Doble Negación $\to$ Distribución). La
pestaña ``Distribución y Leyes'' permite resolver fórmulas de predicados automáticamente.

\subsection{Conjuntos (Equipo Linus)}
Operaciones básicas (unión, intersección, diferencia, complemento, diferencia simétrica) y familia de
conjuntos, con diagramas de Venn interactivos. \textbf{D3.} Aplica la teoría de conjuntos a la lógica
simbólica.

\subsection{Aprender y Progreso}
Vista de recorrido guiado por concepto (definición $\to$ ejemplo $\to$ ejercicio $\to$ solución) con
verificación en vivo, y un store de progreso con \texttt{localStorage} que sugiere temas débiles.

\section{Manual de Uso}
\begin{enumerate}
  \item \textbf{Inicio:} la landing presenta los módulos y la información del curso.
  \item \textbf{Tablas de Verdad:} escribe una fórmula (ej. \texttt{p \^{} q}), inserta conectivos y
    genera la tabla; revisa la clasificación.
  \item \textbf{Inferencias:} introduce premisas y conclusión; el sistema valida la inferencia con trazabilidad.
  \item \textbf{Cuantificadores:} elige $\forall$/$\exists$, define dominio (lista o rango) y
    predicado; usa ``Distribución y Leyes'' para resolver fórmulas paso a paso.
  \item \textbf{Conjuntos:} define conjuntos y aplica operaciones; observa el diagrama de Venn.
  \item \textbf{Aprender:} sigue el recorrido guiado por conectores y leyes.
  \item \textbf{Ejercicios / Progreso:} practica y consulta tu avance y recomendaciones.
\end{enumerate}

\section{Resultados y Aprendizajes}
El proyecto cubrió el 100\% del sílabo con motores funcionales y 175 pruebas. Los aprendizajes por
equipo fueron:
\begin{itemize}
  \item \textbf{Sinergia:} diseño de parsers recursivos y clasificación de fórmulas; comprensión de validez.
  \item \textbf{Los Hijos de Linus:} trazabilidad de inferencias, manejo de errores y formalización de reglas.
  \item \textbf{Modus Innova:} cuantificadores, De Morgan y reescritura de fórmulas; pensamiento algorítmico.
  \item \textbf{Linus:} modelado de conjuntos y visualización; conexión conjuntos--lógica.
\end{itemize}

Estos resultados atañen directamente a \textbf{D1} (lógica proposicional), \textbf{D2} (inferencia) y
\textbf{D3} (conjuntos), demostrando que la construcción de software refuerza la comprensión teórica.

\section{Limitaciones y Trabajo Futuro}
\textbf{Limitaciones:} dominios finitos (no se evalúan cuantificadores sobre infinitos); predicados
libres restringidos a expresiones seguras; la inferencia valida reglas, no demuestra teoremas
arbitrarios.

\textbf{Trabajo futuro:} (i) exportación LaTeX de tablas y demostraciones; (ii) más ejercicios y modo
profesor con calificación; (iii) soporte de lógica de primer orden completa; (iv) tutor adaptativo
basado en el store de progreso.

\section{Conclusiones}
LogiLearn demuestra que la implementación de una plataforma interactiva potencia el aprendizaje de la
lógica simbólica. La metodología por equipos y PRs fomentó trabajo colaborativo y calidad. La
cobertura del 100\% del sílabo con 175 pruebas valida la robustez de los motores. Se recomienda la
adopción de enfoques similares en otros cursos de fundamentos.

\begin{thebibliography}{9}
\bibitem[UNPRG, 2026]{silabo} Universidad Nacional Pedro Ruiz Gallo. (2026). \textit{Sílabo de Lógica Simbólica} (MATG1001, 2026 I).
\bibitem[Copi y Cohen, 2009]{copi} Copi, I. \& Cohen, C. (2009). \textit{Introducción a la lógica} (14.ª ed.). Pearson.
\bibitem[Vue, 2026]{vue} Vue.js. (2026). \textit{Documentación oficial}. \url{https://vuejs.org}
\bibitem[Vite, 2026]{vite} Vite. (2026). \textit{Documentación oficial}. \url{https://vitejs.dev}
\bibitem[Tailwind, 2026]{tw} Tailwind CSS. (2026). \textit{Documentación oficial}. \url{https://tailwindcss.com}
\bibitem[Vitest, 2026]{vitest} Vitest. (2026). \textit{Documentación oficial}. \url{https://vitest.dev}
\end{thebibliography}

\end{document}
